% Standalone section -- paste into the weekly report, or \input it.
% Week 3 working note: whether bifurcation angle, measured the way bifurcation_angle.tex describes,
% separates healthy from diseased cases the way local branch volume does, why it does not, and what
% a fair comparison against the published literature would need. Companion to
% bifurcation_volume.tex, which asks the same disease-separation question of local branch volume and
% finds a signal. Numbers from scripts/bifurcation_volume.py and
% scripts/extract_bifurcation_angles.py, run 2026-09-16 (800 cases); assumes the angle definition
% from bifurcation_angle.tex.

\section*{Bifurcation angle and disease: a well-powered null}

Four independent CCTA cohorts in the published literature report wider bifurcation angle going with
more disease, so the natural next question after measuring angle descriptively
(bifurcation\_angle.tex) was whether the same angle, on the same 800 cases, separates
ImageCAS-X's diseased cases from its healthy ones. It does not, at any of the five named
bifurcations, under any of the tests below, including ones built specifically to be as sensitive as
this cohort's binary disease label allows. That null result turned out to need more explaining than
a positive one would have: whether it is genuinely a gap between angle and disease, or an artifact
of not enough cases, is answered in the first subsection; whether angle is at least hiding a signal
that only shows up alongside volume (bifurcation\_volume.tex) is answered in the second; and what
the literature's own cohorts had that this comparison lacks is worked through last, since a null
result on a two-line label is not the same claim as a null result on angle itself.

\subsection*{LM: well powered for a literature-sized effect, and null}

The left main is the bifurcation the published literature bets on, so it is the fairest place to
ask whether this cohort could have detected the effect those studies report. At the LM, the
per-case pooled standard deviation is $22.4^\circ$ ($n_h = 404$, $n_d = 384$); at that spread, 80\%
power detects an effect of $4.5^\circ$, and the effect actually observed is $2.4^\circ$, below that
floor (\texttt{scripts/extract\_bifurcation\_angles.py}). Reaching 80\% power for the observed
effect would need roughly 1,390 cases per group, well beyond what ImageCAS-X has. But a gradient the
size of the published literature's, $11.8^\circ$ between normal and significantly stenosed arteries
\cite{juan2017bifurcation}, needs only 57 cases per group at our own measured spread, seven times
fewer than we have. So the test run here is not underpowered for the effect the literature claims;
it is underpowered only for the much smaller effect actually present in this cohort's own numbers,
which is a different and more informative statement than "not enough data."

\subsection*{Combining with volume does not recover a signal}

If angle carries some disease information that only becomes visible once the confound of branch
size is accounted for, adding volume as a second predictor should reveal it. It does not: at CRUX,
the bifurcation with the strongest volume signal, angle alone predicts disease status at chance
(5-fold cross-validated AUC $0.488 \pm 0.039$), volume alone reaches $0.576 \pm 0.039$, and angle
added to volume does not improve on volume alone ($0.572 \pm 0.025$). Angle is not merely a weaker
version of the same signal volume carries; on this evidence it carries no signal at all, standing
alone or alongside volume.

\subsection*{Pooling across all five bifurcations does not either}

A single bifurcation's test may simply be unlucky, so the next check pools every named bifurcation
into one per-patient composite: each case's angle at each bifurcation, expressed as a z-score
against that bifurcation's own population, averaged across the bifurcations that case has. That is
the single most powerful angle-only test this label supports, and it is well powered: at this
cohort's own size it detects effects as small as $0.11$ SD at 80\% power. It trends the expected
direction (diseased $+0.038$ SD vs.\ healthy $-0.017$ SD) but does not reach significance
($p = 0.149$, $n_h = 412$, $n_d = 388$). Angle is not underpowered here in the ordinary sense; it is
being asked a question this label cannot answer.

\subsection*{What the literature had that this comparison does not}

None of the four cohorts that motivated this test used a coarse present/absent label, which is the
most direct explanation for why they found an effect and this comparison does not. Temov \& Sun
measured the angle on 196 CCTA cases (mean $79.4^\circ \pm 23.0^\circ$) and found it shifted with
sex and BMI rather than disease alone \cite{temov2016bifurcation}; Juan et al., across 313 cases
stratified by stenosis severity, reported $75.5^\circ$ (normal) rising to $87.3^\circ$ (significant
stenosis) \cite{juan2017bifurcation}; Mansouri et al., in 122 patients with both CCTA and invasive
angiography, correlated angle against the continuous Gensini severity score at $r = 0.684$
($p = 0.0001$), stronger than plaque volume ($0.281$) or calcium score ($0.217$) in the same
patients \cite{mansouri2024bifurcation}; and Tommasino et al.'s CLAP score, following 499 patients,
found an angle above $80^\circ$ the strongest predictor of major adverse cardiac events in that
cohort, HR $4.47$ ($95\%$ CI $3.80$--$6.70$, $p<0.001$), ahead of diabetes and obstructive disease
\cite{tommasino2024clap}. A fifth, smaller cohort reported by a 2026 review of coronary anatomy and
hemodynamics found the same direction in diseased left coronaries ($94^\circ \pm 19.7^\circ$ vs.\
$75.5^\circ \pm 19.8^\circ$, $n=30$, $p=0.02$), while noting that computational studies of the same
relationship disagree on whether angle drives adverse hemodynamics at all \cite{shen2026geometry}.
What every one of the four primary cohorts has and ImageCAS-X's \texttt{Disease} field does not is a
severity grade, a stenosis threshold, or a continuous score, measured at the same site as the angle.
Comparing angle against a patient-level yes/no that can be true because of disease three segments
away dilutes exactly the local effect these studies find; that is a design gap in this comparison,
not evidence that the effect itself is absent.

\subsection*{How far a tighter design alone could close that gap}

Two further findings bound how much of that gap a site-matched, severity-graded redesign could
actually close on its own, without also fixing the measurement itself. Givehchi et al.\ built nine
phantoms of known bifurcation angle, imaged them by CCTA, and measured the same LAD--LCx angle by
two accepted techniques: mean absolute error against the true angle was $2.4^\circ \pm 2.2^\circ$ by
multiplanar reformation and $3.8^\circ \pm 2.9^\circ$ by volume rendering, and applied to 50 real
clinical cases the two techniques disagreed with each other by $12.0^\circ \pm 10.6^\circ$
\cite{givehchi2018phantom}, essentially the same size as Juan et al.'s $11.8^\circ$ disease
gradient. The published angle--disease association sits at the edge of measurement noise even
before a coarse label is introduced, which is why every angle in this thesis comes from one
deterministic, centerline-based definition rather than a manual reformation. Separately, Kwon et
al.\ followed 177 subjects' carotid bifurcation angle over a mean $130.2 \pm 8.1$ months and found
it increased significantly with age alone ($p<0.001$) \cite{kwon2022carotid}; ImageCAS-X carries no
age field at all, so this confound cannot even be checked here, only flagged for any future
angle-outcome comparison on a cohort that does.

\subsection*{What this licenses}

Angle's case for a disease association rests on a site-matched, severity-graded design that this
dataset's binary label cannot provide, not on a claim that the effect is real but small; the tests
above were built specifically to be sensitive to a small effect and still came back null. The ramus
intermedius natural experiment (\texttt{IM} label, Section~\ref{sec:tree}) and the CGPS
baseline-vs-progression design already in the project plan are where that design exists, and where
angle should be tested next, not here.
